% !PRINTED 230
\weberpage{230}

%% [running head, not body text -- commented out 2026-08-07] Page furniture (running
%% head, folio, signature) is excluded from this edition by standing policy.
%% Sixth Section. \hfill \S.\,74.

\begin{center}
\S.\,74.

\gesperrt{Transformation of the equation of the fifth degree.}
\end{center}

\indent These developments shall now be applied in order to carry
out the transformation of the equation of the fifth degree, as we
sketched it in \S.\,54.

\indent A system of values $t_0, t_1, t_2, t_3$ can be determined in
infinitely many ways so that the Bezoutiant $B$ vanishes; in general
the solution of a quadratic equation is required for this. We can,
for example, set
\[
t_2 = 0, \quad t_3 = 0
\]
and determine the ratio $t_0 : t_1$ from the quadratic equation
\[
B_{0,0}\,t_0^2 + 2\,B_{0,1}\,t_0\,t_1 + B_{1,1}\,t_1^2 = 0\,.
\]
These values of $t$ then contain the square root
\[
\sqrt{B_{0,1}^2 - B_{0,0}\,B_{1,1}}\,.
\]
Instead of this, however, one can also make other choices, and one
obtains ever different square roots. In the selection of this square
root there lies something arbitrary and undetermined.

\indent Among these various modes of determination we pick out one
and with it transform the given equation of the fifth degree into a
\gesperrt{principal equation}, that is, into one in which the third
and fourth powers of the unknown do not occur; or we assume, following
F. Klein's precedent, that the equation to be transformed is \gesperrt{from the
outset a principal equation}
\[
% The source's principal-form specialization intentionally has no x^4 or x^3 terms.
\text{(1)} \qquad f(x) = a_0\,x^5 + a_3\,x^2 + a_4\,x + a_5 = 0\,.
\]
But we wish to emphasise expressly that, by \S.\,70, (8), through this
preliminary Tschirnhausen transformation the discriminant of the given
equation is changed only by a factor $\Theta^2$, in which $\Theta$
depends \gesperrt{rationally} on the $t$ employed, and hence also on the square
root used for their determination.

\indent Under the assumption (1) the coefficients of the Bezoutiant,
by the formulas of the preceding paragraph, now become the following:
%CONT>>
